\documentclass[11pt]{article}
\usepackage[T1]{fontenc}
\usepackage{amsmath,amssymb,amsthm}
\usepackage[margin=1in]{geometry}
\usepackage{microtype,array,needspace}
\usepackage[hidelinks]{hyperref}
\hypersetup{pdftitle={A positive semicontinuous density for the Syracuse stationary law},pdfauthor={Lucas Valbuena}}
\newtheorem{theorem}{Theorem}[section]
\newtheorem{proposition}[theorem]{Proposition}
\newtheorem{lemma}[theorem]{Lemma}
\newtheorem{corollary}[theorem]{Corollary}
\theoremstyle{remark}
\newtheorem*{remark}{Remark}
\newcommand{\Zthree}{\mathbb Z_3}
\newcommand{\NN}{\mathbb N_{>0}}
\newcommand{\one}{\mathbf 1}
\newcommand{\HH}{\mathcal H}
\newcommand{\LL}{\mathcal L}
\newcommand{\avg}[1]{\langle #1\rangle}
\DeclareMathOperator{\Law}{Law}
\title{A positive semicontinuous density\\for the Syracuse stationary law}
\author{Lucas Valbuena}
\date{8 October 2026}

\begin{document}
\maketitle

\begin{abstract}
Let \(\mu\) be the stationary probability measure on \(\Zthree\)
for \(x\mapsto(3x+1)/2^a\), with independent geometric exponents
\(\mathbb P(a=j)=2^{-j}\). We construct its greatest nonnegative
lower semicontinuous density with respect to additive Haar probability.
This density vanishes off the units, has a uniform positive lower bound
on the units, and takes a prescribed arithmetic value at every positive
integer. The value is the canonical trace obtained from weighted
first-hit densities, including the cycle factor at periodic targets.
Consequently every unit cylinder of depth \(k\) has mass at least
\(c3^{-k}\), with one constant \(c>0\), and the normalized minimum
cylinder exponent in Tao's 2020 question equals one. At every positive
integer the cylinder densities also converge to the arithmetic trace
in Ces\`aro mean absolute error. The density is infinite on a dense
Haar-null \(G_\delta\) subset of the units, and its integer trace is
unbounded in every unit residue class. The proof combines periodic first-hit
means, a superharmonic bound on inverse layers, and a normalization
argument for an arithmetic lower envelope. It does not settle forward
Collatz convergence or the expected late-occupation bound.
\end{abstract}

\section{The stationary measure and the arithmetic trace}

Write \(\lambda\) for additive Haar probability on \(\Zthree\),
and \(U=\Zthree^\times\), so \(\lambda(U)=2/3\).
The Syracuse stationary measure is the law of
\[
 X=\sum_{j=0}^{\infty}3^j2^{-(a_1+\cdots+a_{j+1})},
 \qquad \mathbb P(a_i=a)=2^{-a}\quad(a\geq1),
\]
where the exponents are independent. This is Tao's random variable
\cite{Tao}. Put
\[
 p_k(b)=\mu\{x:x\equiv b\pmod{3^k}\},\qquad
 \rho_k(x)=3^k p_k(x\bmod3^k).
\]
All Haar densities below use \(\lambda\), rather than Haar
probability normalized on \(U\).

For the arithmetic definitions use the shortcut map
\[
 T(n)=\begin{cases}n/2,&2\mid n,\\(3n+1)/2,&2\nmid n.\end{cases}
\]
Let \(o_r(n)\) count odd entries among
\(n,Tn,\ldots,T^{r-1}n\), and set
\[
 R_r(n)=3^{o_r(n)}2^{-r}.
\]
If \(q\) reaches a positive target \(N\), let \(\tau_N(q)\)
be its first hitting time and define
\[
 w_N(q)=\frac{q}{N}R_{\tau_N(q)}(q).
\]
Set \(w_N(q)=0\) when \(q\) does not reach \(N\).
These weights lie in \([0,1]\). Their natural means have limits,
denoted by \(D_N\). We use the arithmetic density and clock results
in \cite{Basins}; their upstream inputs include the separate
constructions of Mazur and Said Duran described there.

If \(N\) has a positive return, let \(r_N\) be its first positive
return time and put
\[
 G_N=(1-R_{r_N}(N))^{-1}.
\]
Otherwise put \(G_N=1\). Define
\begin{equation}\label{eq:g}
 \delta=\frac12\log\frac43,\qquad
 g(N)=\delta G_N N D_N.
\end{equation}
Here a nonperiodic integer means an integer with no positive return
to itself. It may enter a cycle later.

The construction in this paper starts from the actual function
\(g\), rather than a value assigned to an almost-everywhere density
on the countable set of positive integers. For \(x\in\Zthree\), set
\begin{equation}\label{eq:envelope}
 e_k(x)=\inf_{\substack{n\in\NN\\ n\equiv x\pmod{3^k}}}g(n),
 \qquad E(x)=\sup_{k\geq0}e_k(x).
\end{equation}
The functions \(e_k\) are finite and constant on depth-\(k\)
cylinders. The function \(E\) takes values in \([0,\infty]\).

\begin{theorem}\label{thm:envelope}
The function \(E\) has the following properties.
\begin{enumerate}
\item It is lower semicontinuous, \(\int E\,d\lambda=1\), and
\(\mu=E\lambda\).
\item It vanishes on \(3\Zthree\). There is one constant \(c>0\)
such that \(E(x)\geq c\) for all \(x\in U\).
\item At every positive integer, including every periodic integer,
\(E(n)=g(n)\).
\item If \(h:\Zthree\to[0,\infty]\) is lower semicontinuous and
\(h\lambda=\mu\), then \(h\leq E\) everywhere.
\end{enumerate}
\end{theorem}

The greatest representative uses the extended nonnegative reals.
Corollary~\ref{cor:infinity} below shows that its infinite values
cannot be removed while retaining the greatest-representative property.

\begin{corollary}\label{cor:cylinders}
There is a constant \(c>0\) such that
\[
 p_k(b)\geq c3^{-k}\qquad
 (k\geq1,\ b\in\mathbb Z/3^k\mathbb Z,\ 3\nmid b).
\]
If \(m_k=\min_{3\nmid b}p_k(b)\), then
\[
 c3^{-k}\leq m_k\leq\frac32\,3^{-k},\qquad
 \lim_{k\to\infty}\frac{-\log m_k}{k\log3}=1.
\]
\end{corollary}

Tao introduced this minimum and conjectured the exponent-one
asymptotic in his 2020 discussion \cite{Tao2020}.
Nikpour and Rabbani's preprint \cite{NikpourRabbani} reports an
upper bound of \(1.039979\) for this minimum-atom exponent.
Corollary~\ref{cor:cylinders} gives a constant lower factor.
The proof uses the arithmetic inputs behind \eqref{eq:g}; it is
not a deduction from almost-everywhere integrability alone.
We make no claim of an exhaustive priority search.

\begin{corollary}\label{cor:integer}
For every positive integer \(n\),
\begin{equation}\label{eq:cesaro-absolute}
 \lim_{K\to\infty}\frac1K\sum_{k=0}^{K-1}|\rho_k(n)-g(n)|=0.
\end{equation}
Also \(e_k(n)\uparrow g(n)\). In particular, for each
\(\varepsilon>0\) there is a conductor \(j\) such that
\[
 m\in\NN,\quad m\equiv n\pmod{3^j}
 \quad\Longrightarrow\quad g(m)>g(n)-\varepsilon.
\]
\end{corollary}

The convergence in \eqref{eq:cesaro-absolute} concerns the integer
trace of the cylinder densities. It does not assert convergence of
the individual values \(\rho_k(n)\). None of these conclusions
excludes an additional cycle or a divergent positive orbit.

\begin{corollary}\label{cor:infinity}
The set \(\{x:E(x)=\infty\}\) is a dense \(G_\delta\) subset
of \(U\) and has Haar measure zero. In every residue class prime
to three, \(g\) is unbounded on the positive odd integers in that
class. The function \(E\) is discontinuous at every positive
integer prime to three.
\end{corollary}

\section{First-hit means and a bound on inverse layers}

We record the arithmetic inputs in the normalization used here.
They are supplied by the native first-hit development and the
periodic density extension accompanying \cite{Basins}.

The function \(g\) is finite and nonnegative on \(\NN\), vanishes
on multiples of three, and satisfies one uniform bound
\begin{equation}\label{eq:integer-floor}
 g(n)\geq c_0>0\qquad(3\nmid n).
\end{equation}
Its binary harmonic identity is
\begin{equation}\label{eq:harmonic}
 g(n)=\frac12 g(2n)+\frac32\one_{n\equiv2\ (3)}g((2n-1)/3).
\end{equation}
The native canonical identification gives
\begin{equation}\label{eq:cesaro-input}
 \frac1K\sum_{k<K}\rho_k(n)\longrightarrow g(n)
 \qquad(n\in\NN).
\end{equation}
The cycle factor in \eqref{eq:g} is necessary for both formulas
at periodic targets.

For \(K\geq0\) put
\[
 C_{K,N}(q)=
 \one_{\substack{q\text{ odd},\ q\text{ reaches }N\\
                   o_{\tau_N(q)}(q)\leq K}}
 R_{\tau_N(q)}(q).
\]
The indicator makes this zero when the first hitting time is
undefined. For a function \(p\) of positive period \(M\), write
\(\avg p=M^{-1}\sum_{q=0}^{M-1}p(q)\).

\begin{proposition}[Periodic marked first-hit mean]\label{prop:marked}
Let \(N>0\), and let \(p:\mathbb N\to\mathbb R\) be bounded,
nonnegative, periodic, and zero outside odd integers prime to three.
Then
\begin{equation}\label{eq:marked}
 \lim_{K\to\infty}\frac1K\sum_{q\geq1}C_{K,N}(q)p(q)
       =2\delta N D_N\avg p.
\end{equation}
If \(N\) is nonperiodic, the right side is \(2g(N)\avg p\).
\end{proposition}

\begin{proof}
The periodic source theorem gives the natural mean
\(D_N\avg p\) for \(w_N(q)p(q)\), and partial summation gives
the same logarithmic mean. For every \(\varepsilon>0\), the odd
sources reaching \(N\) for which the bounds
\[
 (1-\varepsilon)\frac{\log q}{2\delta}
 \leq o_{\tau_N(q)}(q)
 \leq(1+\varepsilon)\frac{\log q}{2\delta}
\]
fail form a set of absolute logarithmic density zero. For a bounded
nonnegative mask, its exceptional contribution is bounded by a fixed
multiple of the unmasked exceptional contribution. The same domination
applies to the summable endpoint tail beyond the native exponential
height cutoff. The two-sided clock squeeze therefore converts the
logarithmic source mean to
\[
 \frac1K\sum_{q\geq1}
 \one_{\substack{q\text{ odd},\ q\text{ reaches }N\\o_{\tau_N(q)}(q)\leq K}}
 \frac{w_N(q)p(q)}q
 \longrightarrow 2\delta D_N\avg p.
\]
Multiplication by \(N\) gives \eqref{eq:marked}.
The formal proof first treats masks bounded by one, then rescales a
general bounded mask. It retains the summability and endpoint-tail
hypotheses in the clock squeeze and discharges them by domination
from the actual first-hit estimates.
\end{proof}

Call a nonnegative function \(f:\NN\to\mathbb R\) binary
superharmonic if
\begin{equation}\label{eq:superharmonic}
 \frac12 f(2n)+\frac32\one_{n\equiv2\ (3)}f((2n-1)/3)\leq f(n)
 \qquad(n>0).
\end{equation}

\begin{lemma}[Inverse-layer bound]\label{lem:layer}
If \(f\) is binary superharmonic and \(p\) is bounded,
nonnegative, and satisfies \(p(q)\leq f(q)\) for \(q>0\), then
\begin{equation}\label{eq:layer-cap}
 \sum_{q\geq1}C_{K,N}(q)p(q)\leq(K+1)f(N)
 \qquad(N>0,\ K\geq0).
\end{equation}
\end{lemma}

\begin{proof}
First sum over all pairs \((q,r)\) with \(q\) odd,
\(T^rq=N\), and \(o_r(q)=k\), restricting \(r\) to a finite
range. The weight of a pair is \(R_r(q)p(q)\).
For \(k=0\), the only possible contribution is the zero-time
term, bounded by \(f(N)\). For positive \(k\), split according
to the last edge into \(N\). The even predecessor is \(2N\),
with factor \(1/2\) and unchanged odd count. When it exists, the
odd predecessor is \((2N-1)/3\), with factor \(3/2\) and odd
count reduced by one. Induction on the time cutoff, followed by
\eqref{eq:superharmonic}, bounds each layer by \(f(N)\).
Nonnegative summation removes the time cutoff. Summing over
\(0\leq k\leq K\) gives \((K+1)f(N)\). First hits form a subset
of these pairs, proving the claim. The formal series rearrangements
use the native summability of each unmarked odd-depth layer and
boundedness of \(p\).
\end{proof}

\begin{lemma}[Periodic test bound]\label{lem:periodic-bound}
If \(p\) is bounded, nonnegative, periodic, zero outside odd units,
and \(p(q)\leq g(q)\) for positive \(q\), then \(\avg p\leq1/2\).
\end{lemma}

\begin{proof}
Apply Lemma~\ref{lem:layer} with \(f=g\), which is harmonic.
Use the positive nonperiodic target \(N=5\). Its shortcut orbit
enters \(1,2\), so it never returns to five, and \(g(5)>0\) by
\eqref{eq:integer-floor}. After division by \(K\),
Proposition~\ref{prop:marked} gives
\(2g(5)\avg p\leq g(5)\). Cancel \(g(5)\).
\end{proof}

\section{The lower envelope has mass one}

Every cylinder contains a positive integer, so \(e_k\) in
\eqref{eq:envelope} is finite and continuous. Cylinder refinement
gives \(e_k\leq e_{k+1}\), hence \(E\) is lower semicontinuous.
Also \(E(n)\leq g(n)\) for every positive integer.
Every nonunit cylinder at conductor at least one contains only
integer arguments where \(g\) vanishes. The unit lower bound already
holds at conductor one. Therefore
\begin{equation}\label{eq:envelope-support}
 E=0\text{ on }3\Zthree,\qquad E\geq c_0\text{ on }U.
\end{equation}

\Needspace{8\baselineskip}
\begin{lemma}[Maximality of the arithmetic envelope]\label{lem:maximal}
If \(h:\Zthree\to[0,\infty]\) is lower semicontinuous and
\(h(n)\leq g(n)\) at every positive integer, then \(h\leq E\).
\end{lemma}

\begin{proof}
If \(a<h(x)\), lower semicontinuity supplies a cylinder about
\(x\) on which \(h>a\). Every positive integer in that cylinder
has \(g(n)>a\), so its infimum \(e_k(x)\) is at least \(a\).
Take the supremum over \(a<h(x)\).
\end{proof}

For \(k\geq1\) the integer function
\(p(q)=\one_{q\text{ odd}}e_k(q)\) has period \(2\cdot3^k\)
and satisfies Lemma~\ref{lem:periodic-bound}. Among each pair
\(q,q+3^k\), exactly one integer is odd. Thus
\[
 \avg p=\frac12\int e_k\,d\lambda\leq\frac12.
\]
The same integral bound at \(k=0\) follows by monotonicity.
Consequently
\begin{equation}\label{eq:beta}
 0<\beta:=\int E\,d\lambda\leq1.
\end{equation}
Positivity uses \eqref{eq:envelope-support}; the upper bound uses
monotone convergence.

Define the transfer operator on nonnegative functions by
\[
 (\LL h)(x)=\frac12h(2x)+
 \frac32\one_{x\equiv2\ (3)}h((2x-1)/3).
\]
The second inverse branch is used only on its clopen domain.
This operator preserves lower semicontinuity and Haar integrals.
Indeed, doubling preserves \(\lambda\), while
\(y\mapsto(3y+1)/2\) maps \(\Zthree\) onto the residue class
two with Haar scaling \(1/3\).
At every positive integer \(n\), the inequality \(E\leq g\)
and \eqref{eq:harmonic} give \((\LL E)(n)\leq g(n)\).
Lemma~\ref{lem:maximal} therefore yields \(\LL E\leq E\)
everywhere. Equal finite integrals imply \(\LL E=E\) almost
everywhere.

It follows that \(E\lambda\) is stationary for the binary maps
\(H(x)=x/2\) and \(J(x)=(3x+1)/2\), each with probability
\(1/2\). Iterating the \(H\) part of this equation gives
stationarity for \(H^{a-1}J(x)=(3x+1)/2^a\), with weights
\(2^{-a}\); after \(a\) such terms the remaining mass is
\(\beta2^{-a}\).
Uniqueness of the stationary probability follows from the common
\(3\)-adic contraction factor of these affine maps. Hence
\begin{equation}\label{eq:normalized}
 \beta^{-1}E\lambda=\mu.
\end{equation}

\begin{lemma}[Canonical lower semicontinuous minorants]\label{lem:lsc-bound}
If \(h\) is a nonnegative lower semicontinuous density for \(\mu\),
then \(h(n)\leq g(n)\) for every positive integer \(n\).
\end{lemma}

\begin{proof}
For \(a<h(n)\), all sufficiently fine cylinders about \(n\)
have \(h>a\). Their normalized masses therefore satisfy
\(\rho_k(n)\geq a\) for all sufficiently large \(k\).
Equation~\eqref{eq:cesaro-input} gives \(g(n)\geq a\).
Let \(a\) increase to \(h(n)\).
\end{proof}

Apply this lemma to \(h=E/\beta\), using \eqref{eq:normalized}.
Maximality of the arithmetic envelope gives \(E/\beta\leq E\).
Since \(0<E(1)\leq g(1)<\infty\), evaluation at one implies
\(\beta\geq1\). Together with \eqref{eq:beta}, this proves
\[
 \beta=1,\qquad \mu=E\lambda.
\]
Lemma~\ref{lem:lsc-bound} and Lemma~\ref{lem:maximal} also show
that every nonnegative lower semicontinuous canonical density lies
below \(E\).

\section{Equality at the positive integers}

The restriction \(n\mapsto E(n)\) is finite, nonnegative, and
binary superharmonic. Fix a positive nonperiodic integer \(N\).
For \(j\geq1\), apply Lemma~\ref{lem:layer} with
\(f(n)=E(n)\) and \(p(q)=\one_{q\text{ odd}}e_j(q)\).
Proposition~\ref{prop:marked} and the exact odd-mask mean give
\[
 g(N)\int e_j\,d\lambda\leq E(N).
\]
As \(j\to\infty\), the integrals increase to one. Thus
\(g(N)\leq E(N)\). The reverse inequality is part of the
construction, so \(E(N)=g(N)\) at every positive nonperiodic
integer, without a parity restriction.

To handle periodic points, put \(d(n)=g(n)-E(n)\).
This is nonnegative, is zero at nonperiodic points, and satisfies
\[
 d(n)\leq\frac12d(2n)+
       \frac32\one_{n\equiv2\ (3)}d((2n-1)/3).
\]
A deterministic map has at most one periodic predecessor of any
given point: iterate their common image for one less than a common
multiple of their periods to recover either predecessor.
Consequently, if \(n\) is periodic, all predecessors of \(Tn\)
other than \(n\) have zero defect. It follows that
\[
 d(Tn)\leq R_1(n)d(n),\qquad
 d(T^rn)\leq R_r(n)d(n).
\]
For any positive return \(T^rn=n\), the positive affine
correction implies \(R_r(n)<1\). Thus
\(0\leq d(n)\leq R_r(n)d(n)\) forces \(d(n)=0\).
This completes Theorem~\ref{thm:envelope}.

\section{Cylinder minima and trace consequences}

A unit cylinder of depth \(k\geq1\) lies wholly in \(U\).
Integrating \(E\geq c_0\) over it gives
\(p_k(b)\geq c_0 3^{-k}\).
There are \(2\cdot3^{k-1}\) unit cylinders and their masses sum
to one. Their minimum is therefore at most
\((2\cdot3^{k-1})^{-1}\). These bounds imply
\[
 -\log3+\frac{\log c_0}{k}
 \leq \frac{\log m_k}{k}
 \leq -\log3+\frac{\log(3/2)}{k},
\]
proving Corollary~\ref{cor:cylinders}.

At a positive integer \(n\), lower semicontinuity and
\(E(n)=g(n)\) imply
\[
 \rho_k(n)\geq g(n)-\varepsilon
 \quad\hbox{for all sufficiently large }k.
\]
Together with \eqref{eq:cesaro-input}, this improves the signed
Ces\`aro mean to the absolute-error mean. Explicitly,
\(\min(\rho_k(n),g(n))\to g(n)\), and
\[
 |\rho_k(n)-g(n)|
 =\rho_k(n)+g(n)-2\min(\rho_k(n),g(n)).
\]
Average this identity to obtain \eqref{eq:cesaro-absolute}.
Finally, \(e_k(n)\uparrow E(n)=g(n)\). Choosing one \(e_j(n)\)
above \(g(n)-\varepsilon\), and using its definition as a
cylinder infimum, proves the residue-class assertion in
Corollary~\ref{cor:integer}.

\section{Dense infinite values}

Write \(H(x)=x/2\) and \(J(x)=(3x+1)/2\) as before.
The inequality \(\LL E\leq E\) gives, in the extended
nonnegative reals,
\[
 E(Hx)\geq\tfrac12 E(x),\qquad
 E(Jx)\geq\tfrac32 E(x).
\]
Since \(J(-1)=-1\) and \(E(-1)>0\), the second inequality
forces \(E(-1)=\infty\). The first then gives
\[
 E(-2^{-a})=\infty\qquad(a\geq0).
\]
These points are dense in \(U\): powers of two represent every
unit modulo every power of three. For completeness, elementary
cubing gives
\(4^{3^r}=1+3^{r+1}u_r\) with \(u_r\equiv1\pmod3\).
If a power of two represents a chosen residue modulo \(3^k\),
multiplying it by one of
\(1,4^{3^{k-1}},4^{2\cdot3^{k-1}}\) realizes each of the three
lifts modulo \(3^{k+1}\). Start with the two units modulo three.
Negation and inversion permute the unit residues, giving the
claimed density of \(-2^{-a}\).

Lower semicontinuity makes
\[
 \{E=\infty\}=\bigcap_{M\geq1}\{E>M\}
\]
a \(G_\delta\), and integrability makes it Haar null. Fix a unit
cylinder and a finite threshold \(M\). Choose an infinite-value
point inside the cylinder. Lower semicontinuity supplies a smaller
cylinder on which \(E>M\); it contains a positive odd integer
\(n\), where \(g(n)=E(n)>M\). This proves the unboundedness
assertion. Finally, a finite value of \(E\) at a unit cannot be
a continuity point because every neighborhood contains an infinite
value. Theorem~\ref{thm:envelope} supplies finiteness at each
positive integer, proving Corollary~\ref{cor:infinity}.

\Needspace{12\baselineskip}
\section*{Formalization and verification scope}

The accompanying Lean development defines the same weighted first-hit
density, cycle factor, arithmetic infima, and canonical stationary
measure. The accompanying claim map gives the full theorem-to-source
correspondence. Four principal endpoints in the namespace
\nolinkurl{CollatzCanonical.IntegerStationaryEnvelope} are:
\begin{quote}\small\raggedright
\nolinkurl{canonicalIntegerEnvelope_withDensity}\\
\nolinkurl{canonicalIntegerEnvelope_natCast_eq}\\
\nolinkurl{canonicalMinimumUnitMass_base_three_rate}\\
\nolinkurl{canonicalIntegerEnvelope_infinity_closure}
\end{quote}
The finite probability bound is stated directly for the native
\nolinkurl{syracPMF}, not just for an unrelated stationary model.

All 1,620 modules in the exact internal source closure and the generated
declaration audit were freshly compiled with warnings treated as errors
and individually checked by Lean's official kernel checker: 3,242
successful commands. The audit records 36,651 module-header/constant
rows, 36,539 distinct names, and 3,419 encoded-private rows. The only
permitted axioms are \nolinkurl{propext}, \nolinkurl{Classical.choice},
and \nolinkurl{Quot.sound}. Complete inventories of 53,006 external
compiled objects agree before and after. Those libraries were
fingerprinted and reused; neither an external rebuild nor correspondence
with adjacent source checkouts is claimed. Setup reuse, fresh-process
record verification, and verification after relocation passed.
Relocation copied the internal objects and reused the same external
libraries; it did not recompile the sources. The fresh dependency-download
branch was not exercised.

\section*{Attribution and use of AI}
The stationary Syracuse law, the first-passage framework, and the
minimum-cylinder question are due to Tao. The arithmetic inputs are
those of the companion basin-density development, with its separate
attribution to Mazur and Said Duran. The contribution here is the
arithmetic-envelope construction, its normalization and pointwise
identification, and the stated consequences.

OpenAI Codex was used extensively for mathematical exploration,
Lean development, automated review, and preparation of this manuscript.
Automated checking and review are not external peer review or
endorsement by the authors of the cited works.

\section*{License}
Copyright \textcopyright\ 2026 Lucas Valbuena.
This paper and its source are licensed under
\href{https://creativecommons.org/licenses/by/4.0/}{CC BY 4.0}.
Original accompanying code is licensed under MIT; imported formal
sources retain their own licenses and attribution.

\begin{thebibliography}{9}\small\raggedright
\bibitem{Tao} T.~Tao,
\emph{Almost all orbits of the Collatz map attain almost bounded values},
Forum of Mathematics, Pi \textbf{10} (2022), e12.
\url{https://doi.org/10.1017/fmp.2022.8}.
\bibitem{Tao2020} T.~Tao,
\emph{Equidistribution of Syracuse random variables and density of
Collatz preimages}, 25 January 2020.
\url{https://terrytao.wordpress.com/2020/01/25/equidistribution-of-syracuse-random-variables-and-density-of-collatz-preimages/}.
\bibitem{NikpourRabbani} P.~Nikpour and M.~Rabbani,
\emph{A Certified Density Exponent for the 3x + 1 Problem and the
Smallest Atom of the Syracuse Random Variable}, preprint, 2026.
\url{https://doi.org/10.2139/ssrn.7239062}.
\bibitem{Basins} L.~Valbuena,
\emph{Collatz basin densities and canonical traces from first-passage
stabilization}, 2026, accompanying manuscript and formal development.
\url{https://github.com/x1xhlol/collatz-two-coordinate-obstruction}.
\end{thebibliography}
\end{document}
